\chapter{Ajustements de valeur pour risque de crédit et pour risque de crédit propre}\label{ch:rc:credit-and-debit-valuation-adjustments}

Lorsque le Comité de Bâle a finalisé en juin 2011 ses règles de fonds propres pour le risque de contrepartie,
il en a donné la raison en une phrase~: pendant la crise financière, environ les deux tiers des
pertes imputées au \omterm{def:rc:counterparty-exposure:ccr}{risque de crédit de contrepartie} provenaient de variations de l'\omterm{def:rc:credit-and-debit-valuation-adjustments:cva}{ajustement de valeur
pour risque de crédit}, et seulement un tiers environ de défauts effectifs. Les banques avaient perdu de l'argent
non parce que leurs contreparties avaient fait défaut, mais parce que le prix de marché de la probabilité qu'elles
le fassent avait augmenté, et ce prix entrait dans la fair value de chaque dérivé non garanti.
Ce chapitre évalue cette probabilité~: l'\omterm{def:rc:credit-and-debit-valuation-adjustments:cva}{ajustement de valeur pour risque de crédit} que la banque déduit de la
valeur de ses dérivés pour le défaut de la contrepartie, l'\omterm{def:rc:credit-and-debit-valuation-adjustments:dva}{ajustement de valeur pour risque de crédit propre}
qu'elle y ajoute pour son propre défaut, la manière dont les deux se calculent à partir des expositions du chapitre~17 et des
courbes d'intensité du chapitre~13, puis comment ils sont couverts, comptabilisés et traités en fonds propres.

\section{La CVA comme prix du défaut de la contrepartie}

\begin{definition}[Ajustement de valeur pour risque de crédit]\label{def:rc:credit-and-debit-valuation-adjustments:cva}
L'\emph{ajustement de valeur pour risque de crédit}\index{ajustement de valeur pour risque de crédit} (CVA) d'un \omterm{def:rc:counterparty-exposure:netting}{ensemble
de compensation} (netting set) est la valeur de marché de la perte liée au défaut de la contrepartie~: la valeur des transactions
sans risque de contrepartie moins leur valeur avec ce risque. Avec l'instant de défaut de la contrepartie $\tau_C$,
$R_C$ son taux de recouvrement et l'exposition actualisée,
\begin{align*}
\mathrm{CVA} &= (1-R_C)\,\E\bigl[\mathbf 1_{\tau_C\le T}\,D(0,\tau_C)\max(V_{\tau_C}-C_{\tau_C},0)\bigr]\\
&\approx (1-R_C)\sum_k\mathrm{DEE}(t_k)\bigl(Q_C(t_{k-1})-Q_C(t_k)\bigr),
\end{align*}
où $\mathrm{DEE}(t) = \E[D(0,t)E_t]$ et l'approximation suppose l'exposition et le
défaut indépendants.
\end{definition}

La formule est une somme, sur les dates de défaut, de trois facteurs~: la \omterm{def:rc:reduced-form-credit:pd}{perte en cas de défaut}, la
probabilité que le défaut survienne pendant chaque période (issue de la courbe des swaps de défaut de la contrepartie),
puis l'\omterm{def:rc:counterparty-exposure:profiles}{exposition attendue} actualisée à ce moment. Elle correspond à la valeur d'un swap de défaut conditionnel
dont le notionnel est l'exposition future.

\begin{example}[La CVA de l'ensemble de compensation]\label{ex:rc:credit-and-debit-valuation-adjustments:cva}
La contrepartie de l'\omterm{def:rc:counterparty-exposure:netting}{ensemble de compensation} du chapitre~17 est cotée à 80, 95, 110, 140, 160 et 175
points de base de un à dix ans (BBB, à titre illustratif), avec un taux de recouvrement de 40~\%~: sa \omterm{def:rc:reduced-form-credit:pd}{probabilité
de défaut} à dix ans est de 26.4~\%. Avec les expositions actualisées trajectoire par trajectoire selon les
scénarios de Hull--White, la CVA de l'\omterm{def:rc:counterparty-exposure:netting}{ensemble de compensation} s'élève à 891\,342~USD. Isolément, la CVA du swap vaut 288\,250 et
celle du swap de devises 768\,334~: la compensation fait économiser 16~\%. Sur l'\omterm{def:rc:reduced-form-credit:risky}{annuité risquée}
de 7.45 du swap de défaut à dix ans, la CVA vaut 12.0 points de base par an sur 100 millions d'USD
(voir la \cref{fig:rc:credit-and-debit-valuation-adjustments:contrib}).
\end{example}

\begin{omfigure}
\begin{tikzpicture}
\begin{axis}[width=11.5cm,height=5.2cm,ybar,bar width=14pt,
  xlabel={année de défaut},ylabel={contribution à la CVA (millions d'USD)},
  xtick={0.5,1.5,2.5,3.5,4.5,5.5,6.5,7.5,8.5,9.5},xticklabels={1,2,3,4,5,6,7,8,9,10},
  tick label style={font=\small},label style={font=\small},
  xmin=0,xmax=10,ymin=0,ymax=0.14,grid=major,grid style={gray!25},
  yticklabel style={/pgf/number format/fixed,/pgf/number format/precision=2},
  table/col sep=comma]
\addplot[fill=omDef!70,draw=omDef] table[x=year,y=cva]{figdata/rates-credit-risk/18-credit-and-debit-valuation-adjustments/contrib.csv};
\end{axis}
\end{tikzpicture}

\omcaption{La CVA de l'\omterm{def:rc:counterparty-exposure:netting}{ensemble de compensation} par année de défaut de la contrepartie~: le produit de
l'\omterm{def:rc:counterparty-exposure:profiles}{exposition attendue} actualisée (croissante, puis stable) et de la \omterm{def:rc:reduced-form-credit:pd}{probabilité de défaut} de l'année
(croissante avec la courbe d'intensité, puis actualisée par la survie). Les années quatre à six contribuent le plus.
Données~: tutoriel du chapitre.}\label{fig:rc:credit-and-debit-valuation-adjustments:contrib}
\end{omfigure}

\omcode{code/firm/cva/firm_cva.py}{42}{48}{CVA à partir d'un profil d'exposition attendue actualisé et d'une courbe d'intensité, unilatérale ou en premier défaut.}\label{lst:rc:credit-and-debit-valuation-adjustments:cva}

\section{DVA et vision bilatérale}

\begin{definition}[Ajustement de valeur pour risque de crédit propre]\label{def:rc:credit-and-debit-valuation-adjustments:dva}
L'\emph{ajustement de valeur pour risque de crédit propre}\index{ajustement de valeur pour risque de crédit propre} (DVA) est la même
grandeur vue depuis la contrepartie~: la valeur de marché du gain pour la banque lié à son
propre défaut, lorsqu'elle cesserait de payer ce qu'elle doit, calculée à partir de l'exposition négative actualisée
et de la courbe d'intensité et du taux de recouvrement propres à la banque.
\end{definition}

\begin{definition}[CVA bilatérale]\label{def:rc:credit-and-debit-valuation-adjustments:bcva}
La \emph{CVA bilatérale}\index{CVA bilatérale} est l'ajustement $\mathrm{CVA}-\mathrm{DVA}$ lorsque chacun est
calculé en premier défaut~: le défaut de la contrepartie ne coûte à la banque que si la banque
n'a pas fait défaut auparavant (poids $Q_B(t)$ dans la somme de la CVA), et réciproquement. Elle fait que les deux
parties s'accordent sur le prix~: la CVA de l'une est la DVA de l'autre.
\end{definition}

\begin{example}[Les deux côtés]\label{ex:rc:credit-and-debit-valuation-adjustments:bcva}
La banque elle-même est cotée à 40 à 90 points de base (single-A, à titre illustratif). Son \omterm{def:rc:counterparty-exposure:netting}{ensemble de compensation}
dérive en faveur de la contrepartie (chapitre~17), si bien que sa DVA de 1\,055\,437~USD dépasse sa CVA.
En premier défaut, la CVA vaut 832\,416 et la DVA 896\,328~: l'ajustement bilatéral est de
$-\text{USD}~63\,912$, un petit \emph{bénéfice} net pour la banque. Avec la CSA à seuil nul du chapitre~17,
la CVA tombe à 122\,340~USD et la DVA à 70\,490.
\end{example}

\begin{remark}[Une banque peut-elle monétiser son propre défaut~?]\label{rem:rc:credit-and-debit-valuation-adjustments:dva}
La DVA est une valeur réelle pour les créanciers de la banque, non pour ses actionnaires~: elle ne se réalise que
en faisant défaut. Elle ne peut pas être couverte en achetant une protection sur soi-même~; les banques approchent
la couverture en vendant une protection sur des banques similaires, ce qui ajoute un risque de corrélation. Sa croissance lorsque
le crédit de la banque se dégrade est la raison pour laquelle les régulateurs l'écartent des fonds propres.
\end{remark}

\begin{definition}[Montant de résiliation]\label{def:rc:credit-and-debit-valuation-adjustments:closeout}
Le \emph{montant de résiliation}\index{montant de résiliation} (close-out amount) est le montant que la partie survivante
détermine, selon la convention-cadre, pour l'\omterm{def:rc:counterparty-exposure:netting}{ensemble de compensation} résilié~: en gros le coût de
remplacement des transactions, qui peut inclure la vision qu'a la contrepartie de remplacement du
crédit propre du survivant (résiliation risquée) ou non (résiliation sans risque). Ce choix modifie la formule de la CVA~: avec une résiliation risquée, la
DVA de la partie survivante survit au défaut de la contrepartie.
\end{definition}

\section{Calcul de la CVA}

\begin{method}[Calcul de la CVA dans un moteur de simulation]\label{met:rc:credit-and-debit-valuation-adjustments:engine}
(1) Simuler les expositions et les déflateurs sur des scénarios communs (chapitre~17). (2) Calculer les
expositions positives et négatives attendues actualisées de chaque \omterm{def:rc:counterparty-exposure:netting}{ensemble de compensation}. (3) Bootstrapper
la courbe d'intensité de chaque contrepartie à partir de ses cotations de swaps de défaut, ou d'une courbe de substitution pour un nom
sans cotation (chapitre~20). (4) Sommer sur la grille. (5) Pour le \omterm{def:rc:counterparty-exposure:wwr}{risque de corrélation défavorable}, simuler
l'intensité de défaut avec les facteurs de marché et moyenner la perte trajectorielle.
\end{method}

\begin{example}[CVA en risque de corrélation défavorable]\label{ex:rc:credit-and-debit-valuation-adjustments:wwr}
Lier l'intensité de défaut de la contrepartie à l'euro~:
\[
\lambda_t = \lambda^0_t\,\bigl(X_t/F_X(0,t)\bigr)^{-b},
\]
normalisée de sorte
que la courbe de survie moyenne reste celle du marché, comme au chapitre~17. La CVA du swap de devises
vaut 766\,532~USD sans dépendance ($b = 0$), 1\,108\,459 pour $b = 2$ et 1\,521\,524 pour
$b = 5$~: le \omterm{def:rc:counterparty-exposure:wwr}{risque de corrélation défavorable} (wrong-way risk) double la CVA (voir la \cref{fig:rc:credit-and-debit-valuation-adjustments:wwr}).
Pour une dépendance très forte, la CVA diminue légèrement~: la contrepartie fait alors défaut tôt
sur les trajectoires où l'euro baisse en premier, avant que l'exposition n'ait augmenté.
\end{example}

\begin{omfigure}
\begin{tikzpicture}
\begin{axis}[width=11.5cm,height=5.2cm,
  xlabel={dépendance $b$ de l'intensité de défaut à l'euro},ylabel={CVA (millions d'USD)},
  tick label style={font=\small},label style={font=\small},
  xmin=0,xmax=10,ymin=0,ymax=1.8,grid=major,grid style={gray!25},
  table/col sep=comma]
\addplot[omThm,very thick,mark=*,mark size=1.5pt] table[x=b,y=cva]{figdata/rates-credit-risk/18-credit-and-debit-valuation-adjustments/wwr.csv};
\end{axis}
\end{tikzpicture}

\omcaption{CVA du swap de devises lorsque l'intensité de défaut de la contrepartie dépend plus fortement
de l'euro, la courbe de survie moyenne étant maintenue à celle du marché. Le \omterm{def:rc:counterparty-exposure:wwr}{risque de corrélation défavorable} double la
CVA~; la courbe s'infléchit lorsque les défauts se concentrent si tôt que l'exposition reste faible.
Données~: le tutoriel du chapitre.}\label{fig:rc:credit-and-debit-valuation-adjustments:wwr}
\end{omfigure}

\section{Sensibilités et couverture}

La CVA varie avec tous les facteurs dont dépend l'exposition (taux, change, volatilités) et avec le crédit de la
contrepartie. Un desk de CVA couvre le premier ensemble avec les instruments sous-jacents, et
l'autre par des swaps de défaut sur la contrepartie, ou sur un indice lorsque le nom ne
se négocie pas.

\begin{example}[CS01 par compartiments]\label{ex:rc:credit-and-debit-valuation-adjustments:cs01}
En relevant chacun des six spreads cotés de la contrepartie d'un point de base, la courbe étant
de nouveau bootstrappée, modifie la CVA de $-235$, $-165$, $-141$, $-22$, $+125$ et $+4\,693$ dollars
(voir la \cref{fig:rc:credit-and-debit-valuation-adjustments:cs01})~: une hausse parallèle coûte
4\,254~USD par point de base. Presque tout se concentre sur la cotation à dix ans~; relever une cotation courte
avec les longues fixées abaisse les intensités à terme ultérieures, si bien que les compartiments courts sont négatifs.
La couverture représente environ 5.7~millions de dollars de protection à dix ans (le \omterm{def:rc:reduced-form-credit:cs01}{CS01} divisé par
l'\omterm{def:rc:reduced-form-credit:risky}{annuité risquée} du swap de défaut et d'un point de base).
\end{example}

\begin{omfigure}
\begin{tikzpicture}
\begin{axis}[width=11.5cm,height=5.2cm,ybar,bar width=18pt,
  symbolic x coords={1y,2y,3y,5y,7y,10y},xtick=data,xticklabels={1A,2A,3A,5A,7A,10A},
  xlabel={spread au pair de la contrepartie, choqué},ylabel={variation de la CVA (USD par pb)},
  tick label style={font=\small},label style={font=\small},
  ymin=-1200,ymax=5600,grid=major,grid style={gray!25},
  nodes near coords,every node near coord/.append style={font=\footnotesize,/pgf/number format/fixed,/pgf/number format/precision=0},
  table/col sep=comma]
\addplot[fill=omThm!60,draw=omThm] table[x=tenor,y=cs01]{figdata/rates-credit-risk/18-credit-and-debit-valuation-adjustments/cs01.csv};
\end{axis}
\end{tikzpicture}

\omcaption{\omterm{def:rc:reduced-form-credit:cs01}{CS01} par compartiments de la CVA de l'\omterm{def:rc:counterparty-exposure:netting}{ensemble de compensation}. La cotation à dix ans porte presque toute la
sensibilité~; les cotations plus courtes ont de faibles sensibilités négatives via le bootstrap. Données~: le
tutoriel du chapitre.}\label{fig:rc:credit-and-debit-valuation-adjustments:cs01}
\end{omfigure}

\section{Comptabilité et réglementation}

\begin{definition}[XVA]\label{def:rc:credit-and-debit-valuation-adjustments:xva}
L'acronyme \emph{XVA}\index{XVA} désigne la famille des ajustements de valorisation qui transforment la
valeur sans risque d'un dérivé en prix facturé ou déclaré par une banque~: la CVA et la DVA pour le défaut, ainsi que les
ajustements de financement, de marge et de capital du chapitre~19.
\end{definition}

Les normes comptables évaluent les dérivés à la fair value, un prix de sortie~: la norme IFRS~13, publiée en mai
2011, la définit avec les hypothèses que retiendraient les intervenants de marché, y compris celles relatives au risque.
Le crédit de la contrepartie est un tel risque, de sorte que la CVA fait partie de la fair value publiée, et il en va de même pour le
crédit propre de la banque, sous forme de DVA~: la fair value d'un passif reflète son risque de non-exécution, lequel
comprend le risque de crédit propre de l'entité (paragraphe 42). Le troisième trimestre 2011 de JPMorgan, par exemple,
a enregistré un gain avant impôt de 1.9~milliard de dollars au titre de la DVA, issu de l'élargissement de ses propres spreads de crédit. La réglementation adopte une autre position~: les pertes de CVA exigent du capital (la
charge introduite par le Comité de Bâle après la crise, chapitre~19), et la DVA, qui croît lorsque
la banque se dégrade, est déduite des fonds propres durs.

\begin{dated}{2026-09}{CVA et DVA dans les fonds propres}\label{dat:rc:credit-and-debit-valuation-adjustments:dva}
Le Comité de Bâle a indiqué en juin 2011 qu'environ deux tiers des pertes de la crise liées au
risque de contrepartie étaient des pertes de CVA. Bâle~III exige que les banques retirent des
fonds propres de base de catégorie~1 les gains et pertes latents liés aux variations de leur propre risque de crédit sur
les passifs évalués à la fair value~; sa règle finale de juillet 2012 a étendu cette règle à toutes les DVA sur dérivés,
avec une application par paliers, de 20~\% en 2014 à une déduction complète à compter du 1~janvier 2018.
\end{dated}

\begin{omfigure}
\begin{tikzpicture}
\begin{axis}[width=11.5cm,height=5.2cm,
  xlabel={élargissement des spreads propres de la banque (pb)},ylabel={DVA (millions d'USD)},
  tick label style={font=\small},label style={font=\small},
  xmin=0,xmax=100,ymin=0,ymax=2.2,grid=major,grid style={gray!25},
  table/col sep=comma]
\addplot[omDef,very thick,mark=*,mark size=1.5pt] table[x=shift,y=dva]{figdata/rates-credit-risk/18-credit-and-debit-valuation-adjustments/dva.csv};
\end{axis}
\end{tikzpicture}

\omcaption{DVA de l'\omterm{def:rc:counterparty-exposure:netting}{ensemble de compensation} lorsque les spreads des swaps de défaut propres de la banque s'élargissent en parallèle~:
plus le crédit de la banque se dégrade, plus le gain comptable est élevé. Données~: le
tutoriel du chapitre.}\label{fig:rc:credit-and-debit-valuation-adjustments:dva}
\end{omfigure}

\section{Tutoriel~: CVA d'un ensemble de compensation}\label{tut:rc:credit-and-debit-valuation-adjustments:tutorial}

\begin{tutorial}
\textbf{Objectif.} Calculer la CVA et la DVA de l'\omterm{def:rc:counterparty-exposure:netting}{ensemble de compensation} du chapitre~17, unilatérales et
first-to-default, avec et sans CSA, leurs sensibilités au crédit, et l'effet du
\omterm{def:rc:counterparty-exposure:wwr}{risque de corrélation défavorable} (wrong-way risk). \textbf{État final~:} les
\cref{ex:rc:credit-and-debit-valuation-adjustments:cva,ex:rc:credit-and-debit-valuation-adjustments:bcva,ex:rc:credit-and-debit-valuation-adjustments:wwr,ex:rc:credit-and-debit-valuation-adjustments:cs01}
et les quatre graphiques.
\begin{tutsteps}
\item \textbf{Entrées}~: scénarios et valeurs du chapitre~17~; \texttt{deflators(sc)}~;
\texttt{bootstrap} pour les deux courbes de crédit.
\item \textbf{Ajustements}~: \texttt{adjustments()} et \texttt{adjustments(csa=True)}.
\item \textbf{Risque}~: \texttt{cs01()}~; \texttt{wrong\_way()}.
\item \textbf{DVA}~: \texttt{dva\_shift()}~; \texttt{fig\_rc\_cva.py} écrit les graphiques.
\end{tutsteps}
\textbf{Pour aller plus loin.} Remplacer la courbe de la contrepartie par une courbe plate de 175 points de base et
comparer~; calculer la CVA avec l'EE non actualisée multipliée par l'obligation zéro-coupon, et mesurer
l'erreur commise en ignorant la corrélation entre les taux et l'exposition.
\end{tutorial}

\section{Construction~: le moteur de CVA}\label{bld:rc:credit-and-debit-valuation-adjustments:cva}

\begin{build}
\bfield{Rôle} La CVA et la DVA de la banque~: ajustements de juste valeur de chaque \omterm{def:rc:counterparty-exposure:netting}{ensemble de compensation}, leurs
sensibilités au crédit pour le desk de couverture, et les entrées de la valorisation du chapitre~20.
\bfield{Interface} \texttt{deflators(sc)}~; \texttt{discounted\_profiles(E, D)}~;
\texttt{cva(dee, times, cpty, recovery, own=None)}~; \texttt{dva}~; \texttt{cva\_pathwise(E, D,
times, hazards)}~; \texttt{cs01\_buckets}~; \texttt{running\_spread}.
\bfield{Règles} Exposition aux milieux des périodes~; indépendance sauf si une intensité est simulée~;
taux de recouvrement de 40~\% par défaut~; courbes issues de \texttt{firm\_cdscurve}.
\bfield{Tests d'acceptation} \texttt{code/firm/cva/tests/}~: une exposition et une intensité plates donnent la
forme fermée~; le moteur par trajectoires avec une intensité déterministe est égal au moteur par profils~;
la CVA du premier défaut est inférieure à la CVA unilatérale~; les \omterm{def:rc:reduced-form-credit:cs01}{CS01} par compartiments s'additionnent au \omterm{def:rc:reduced-form-credit:cs01}{CS01} parallèle.
\bfield{Pour aller plus loin} Taux de recouvrement stochastique~; close-out risqué~; grecques de la CVA par
différentiation algorithmique adjointe~; CVA incrémentale (chapitre~20).
\end{build}

\begin{omsources}
\item Comité de Bâle sur le contrôle bancaire, communiqué de presse du 1 juin 2011 sur le traitement en capital
du risque de crédit de contrepartie bilatéral.
\item Comité de Bâle sur le contrôle bancaire, \emph{Basel~III}, révisé en juin 2011,
paragraphe~75~; communiqué de presse du 25 juillet 2012 sur les ajustements de valeur des
passifs dérivés.
\item Fondation IFRS, IFRS~13 \emph{Fair Value Measurement}, 2011.
\end{omsources}

\section{Exercices}

\begin{exercise}[$\star$]\label{exo:rc:credit-and-debit-valuation-adjustments:1}
Un \omterm{def:rc:counterparty-exposure:netting}{ensemble de compensation} a une EE actualisée plate de 2~millions de dollars pendant cinq ans~; l'intensité
de défaut de la contrepartie est plate à 2~\% et son taux de recouvrement de 40~\%. Donner la CVA.
\end{exercise}

\begin{exercise}[$\star$]\label{exo:rc:credit-and-debit-valuation-adjustments:2}
Estimer la CVA du chapitre comme l'EPE multipliée par le spread à dix ans et par l'\omterm{def:rc:reduced-form-credit:risky}{annuité risquée}, et
comparer.
\end{exercise}

\begin{exercise}[$\star$]\label{exo:rc:credit-and-debit-valuation-adjustments:3}
Pourquoi l'ajustement bilatéral de l'\cref{ex:rc:credit-and-debit-valuation-adjustments:bcva}
est-il négatif, et que voit la contrepartie~?
\end{exercise}

\begin{exercise}[$\star\star$]\label{exo:rc:credit-and-debit-valuation-adjustments:4}
Pourquoi la CVA de l'\omterm{def:rc:counterparty-exposure:netting}{ensemble de compensation} est-elle inférieure à la somme des CVA autonomes~?
\end{exercise}

\begin{exercise}[$\star\star$]\label{exo:rc:credit-and-debit-valuation-adjustments:5}
Quelle protection à dix ans couvre la sensibilité au crédit de la CVA, et que
laisse ouvert la couverture~?
\end{exercise}

\begin{exercise}[$\star\star$]\label{exo:rc:credit-and-debit-valuation-adjustments:6}
Pourquoi le CSA réduit-il la CVA de 86~\% mais la DVA de 93~\%~?
\end{exercise}

\begin{exercise}[$\star\star\star$]\label{exo:rc:credit-and-debit-valuation-adjustments:7}
\emph{Programmation.} Calculer la CVA en cas de corrélation défavorable pour $b = 0$, 2 et 5, et expliquer le
retournement de la courbe pour les grandes valeurs de $b$.
\end{exercise}

\begin{exercise}[$\star\star\star$]\label{exo:rc:credit-and-debit-valuation-adjustments:8}
\emph{Trouver la faille.} «~Le gain de DVA de notre trimestre prouve que notre activité de produits dérivés vaut
plus que ne le pensait le marché.~»
\end{exercise}

\section{Problème~: le gain de devenir plus risqué}

\begin{problem}[{Problème du week-end --- crédit propre et compte de résultat}]\label{pb:rc:credit-and-debit-valuation-adjustments:1}
Dans un trimestre de tensions, les spreads des swaps de défaut propres à la banque s'élargissent de 50 points de base à
toutes les échéances, sans que rien d'autre ne bouge. Considérer l'\omterm{def:rc:counterparty-exposure:netting}{ensemble de compensation} du chapitre~17 et la courbe de la banque de
l'\cref{ex:rc:credit-and-debit-valuation-adjustments:bcva}.

\medskip\noindent\textbf{Partie I --- Le gain.}
\begin{enumerate}
\item Donner la DVA unilatérale avant et après.
\item Donner le gain, et indiquer où il apparaît dans les comptes.
\item Pourquoi la DVA augmente-t-elle lorsque le crédit de la banque se dégrade~?
\item Rapporter le gain à un portefeuille de produits dérivés 500 fois plus grand.
\item Qui est moins bien loti au cours du même trimestre~?
\end{enumerate}

\medskip\noindent\textbf{Partie II --- Fonds propres.}
\begin{enumerate}[resume]
\item Que fait Bâle~III du gain~?
\item Quand la déduction a-t-elle été pleinement introduite~?
\item Pourquoi les régulateurs traitent-ils la DVA différemment de la CVA~?
\item Que deviendrait le ratio de fonds propres de base de catégorie~1 de la banque si le gain était comptabilisé~?
\item Comment les investisseurs liraient-ils le gain~?
\end{enumerate}

\medskip\noindent\textbf{Partie III --- Le couvrir.}
\begin{enumerate}[resume]
\item Pourquoi la banque ne peut-elle pas acheter de protection sur elle-même~?
\item Quel produit de couverture par substitution pourrait-elle utiliser, et quel risque ajoute-t-il~?
\item La banque doit-elle couvrir la DVA~?
\item Comment un CSA modifie-t-il la DVA et le gain~?
\item Que voit la contrepartie dans sa CVA~?
\end{enumerate}

\medskip\noindent\textbf{Partie IV --- Jugement.}
\begin{enumerate}[resume]
\item La DVA a-t-elle une valeur~? Pour qui~?
\item Comment le P\&L d'un trader doit-il traiter la DVA~?
\item Pourquoi l'ajustement bilatéral est-il celui sur lequel les deux parties peuvent s'accorder~?
\item Énoncer le \emph{résultat nommé}~: le gain de DVA lié à un élargissement de 50 points de base sur le
\omterm{def:rc:counterparty-exposure:netting}{ensemble de compensation}.
\item En une phrase~: pourquoi la DVA est-elle déduite des fonds propres~?
\end{enumerate}
\end{problem}

\section{Questions d'entretien}

\begin{interviewq}[$\star$ \iqroles{trader, risque}]\label{iq:rc:credit-and-debit-valuation-adjustments:1}
Écrivez la formule de la CVA et expliquez chacun de ses termes.
\end{interviewq}

\begin{interviewq}[$\star\star$ \iqroles{chercheur}]\label{iq:rc:credit-and-debit-valuation-adjustments:2}
Qu'est-ce que la DVA, et pourquoi suscite-t-elle la controverse~?
\end{interviewq}

\begin{interviewq}[$\star\star$ \iqroles{trader}]\label{iq:rc:credit-and-debit-valuation-adjustments:3}
Comment couvririez-vous la CVA d'un portefeuille de swaps non garantis avec une entreprise notée BBB~?
\end{interviewq}

\begin{interviewq}[$\star\star$ \iqroles{développeur, chercheur}]\label{iq:rc:credit-and-debit-valuation-adjustments:4}
Comment intégrez-vous le \omterm{def:rc:counterparty-exposure:wwr}{risque de corrélation défavorable} dans un moteur de CVA~?
\end{interviewq}

\begin{interviewq}[$\star\star\star$ \iqroles{risque, banque}]\label{iq:rc:credit-and-debit-valuation-adjustments:5}
Pourquoi la CVA a-t-elle causé davantage de pertes que les défauts pendant la crise~?
\end{interviewq}

\begin{interviewq}[$\star\star\star$ \iqroles{chercheur}]\label{iq:rc:credit-and-debit-valuation-adjustments:6}
Close-out sans risque ou close-out risqué~: quelle est la différence et pourquoi est-elle importante~?
\end{interviewq}
